\section{Conclusion et perspectives}
\label{sec:conclusion}

\subsection{Synthèse}

Ce rapport a présenté un pipeline ML en deux phases pour la classification IRM
de tumeurs cérébrales. La Phase~1 a utilisé trois métriques de séparabilité (FDR,
MI, DBI) pour sélectionner les paramètres optimaux de quatre familles de descripteurs
artisanaux sur \NPhaseOneConfigs{} configurations, sans entraînement de classifieur.
La Phase~2 a benchmarké \NModels{} classifieurs sur \NFeatureSets{} jeux de
caractéristiques (\NRuns{} pipelines évalués), avec une validation statistique complète.

Les principaux résultats sont les suivants :
\begin{itemize}
  \item \textbf{Meilleur pipeline :} Full(opt) + SVM --- précision \BestAcc{},
        F1-macro \BestFOne, $\kappa = \BestKappa$.
  \item \textbf{La fusion améliore les familles isolées :} + \FusionGain{} points
        de pourcentage par rapport à la référence HOG seul.
  \item \textbf{Robustesse statistique :} variance CV $\leq 0{,}01$, IC bootstrap
        [\BestFOneCILow, \BestFOneCIHigh], F1 moyen par graine \SeedFOneMean.
  \item \textbf{Significativité du SVM :} significativement meilleur que RF et
        XGBoost sur Full(opt) (McNemar, $p < 0{,}02$).
\end{itemize}

\subsection{Perspectives}

\begin{enumerate}
  \item \textbf{Comparaison avec l'apprentissage profond :} évaluer des CNN
        (EfficientNet, ResNet) et des Vision Transformers sur le même jeu de données
        et protocole pour une comparaison directe.
  \item \textbf{Systèmes hybrides ML/DL :} combiner des caractéristiques artisanaux
        optimisées avec des embeddings appris en profondeur pour un pool de
        caractéristiques complémentaire.
  \item \textbf{IA explicable :} intégrer Grad-CAM ou SHAP pour fournir des
        explications spatiales aux côtés de la décision du classifieur, renforçant
        la confiance clinique.
  \item \textbf{Validation multi-centres :} évaluer le pipeline sur des jeux de
        données hétérogènes en termes d'appareils pour mesurer la généralisabilité clinique.
\end{enumerate}

L'apprentissage automatique classique avec des descripteurs artisanaux optimisés
reste une approche compétitive et transparente pour la classification IRM de
tumeurs cérébrales, en particulier dans des scénarios à données limitées et à
forte exigence d'interprétabilité.
